\documentclass{article}
\usepackage[utf8]{inputenc}
\usepackage{hyperref}
\usepackage{listings}
\usepackage{xcolor}
\usepackage{enumitem}

\lstset{
    basicstyle=\ttfamily\small,
    breaklines=true,
    frame=single,
    backgroundcolor=\color{gray!5},
}

\title{Display disk space usage for all directories with hierarchical breakdown}
\author{marcos}
\date{2026-04-21}

\begin{document}

\maketitle

\begin{abstract}
Shows disk space usage for directories up to 2 levels deep in the current directory tree, sorted by size in descending order with human-readable format, providing a hierarchical breakdown of disk consumption.
\end{abstract}

\section*{Language}
bash

\section*{Category}
disk-usage

\section*{Command}
\begin{lstlisting}
du -h --max-depth=2 . | sort -hr
\end{lstlisting}

\section*{Explanation}
The command uses 'du -h' to calculate disk usage in human-readable format. The '--max-depth=2' flag limits output to directories up to 2 levels deep from the current directory, avoiding excessive detail while still showing the hierarchy. The '.' specifies the starting directory. The output is piped to 'sort -hr' to sort in descending order by human-readable sizes, making it easy to identify which directories consume the most space at each level.

\section*{Tags}
du, disk-usage, directories, storage, hierarchical, sort, system-administration

\section*{Dependencies}
coreutils

\section*{Arguments}
\begin{enumerate}
  \item \textbf{PATH} (Optional): Root directory to analyze (default: current directory) \\ Default: .
  \item \textbf{DEPTH} (Optional): Maximum directory depth to traverse (default: 2) \\ Default: 2
\end{enumerate}

\section*{Examples}
\begin{enumerate}
  \item \texttt{du -h --max-depth=2 . | sort -hr} - Show disk usage up to 2 levels deep from current directory
  \item \texttt{du -h --max-depth=1 . | sort -hr} - Show only immediate subdirectories (1 level deep)
  \item \texttt{du -h --max-depth=3 /home | sort -hr} - Show disk usage in /home up to 3 levels deep
  \item \texttt{du -h --max-depth=2 . | sort -hr | head -20} - Show top 20 directories by size (2 levels deep)
  \item \texttt{du -ah --max-depth=2 . | sort -hr | head -15} - Include files in the listing, show top 15 items
\end{enumerate}

\section*{Output}
\begin{lstlisting}
1.2T	.
847G	./data
512G	./data/cache
256G	./data/backups
244G	./media
128G	./media/videos
116G	./media/archives
156G	./documents
96G	./documents/projects
60G	./documents/reports
\end{lstlisting}

\section*{Notes}
\begin{itemize}
  \item The --max-depth flag controls recursion depth: 0 shows only the starting directory, 1 shows immediate children, etc.
  \item The -h flag makes output human-readable (K, M, G, T); use -B 1M for sizes in fixed block units
  \item The -s flag is not needed here as du summarizes by default
  \item Hidden directories (those starting with .) are included automatically
  \item Each line shows cumulative size for that directory and all its contents
  \item Symbolic links are not followed by default; add -L to follow symlinks
  \item On filesystems with many nested directories, increasing --max-depth may significantly increase runtime
\end{itemize}

\section*{Warnings}
\begin{itemize}
  \item Results may be incomplete if you lack read permissions on some directories
  \item On network filesystems (NFS, SMB), this command can be slow due to network latency
  \item Very deep directory hierarchies combined with high --max-depth values may cause excessive output
  \item Sparse files may report misleading sizes on some filesystems
\end{itemize}

\section*{See Also}
\begin{itemize}
  \item disk-space-sort-largest-directories
\end{itemize}

\section*{Status}
reviewed

\section*{Safety}
safe

\section*{Shell}
bash

\section*{Platforms}
linux, gnu-linux, freebsd, openbsd, netbsd

\section*{Created At}
2026-04-21

\section*{Updated At}
2026-04-21

\end{document}
